% Lösungen Einheit 1 von 5 – Körper erkennen, Netze, Schrägbilder (zusammenbau v0.1)
\einheitenkopf{Einheit 1 von 5 · Körper erkennen, Netze, Schrägbilder}

%% TODO Lösung der Erklärzeile 9g) fehlt in der Bank
\erg{9}{a) Prisma \quad b) Quader \quad c) Dreiecksprisma (Prisma mit Dreiecken als Grund- und Deckfläche) \quad d) Zylinder \quad e) Kegel \quad f) Würfel \quad g) \ldots \quad h) $8$; $12$; $6$ (Ecken; Kanten; Flächen) \quad i) ja \quad j) D – zwischen B und D liegt genau eine Fläche \quad k) sechs Rechtecke, je zwei gleich: $4 \times 2$, $4 \times 1$ und $2 \times 1$ (in cm), so aneinander, dass sie zu einem Quader falten \quad l) Vorderseite $6$ cm $\times$ $3$ cm in wahrer Größe; Tiefe $4$ cm schräg als $2$ cm (Kästchendiagonalen); verdeckte Kanten gestrichelt \quad m) Diagonalen der Grundfläche zeichnen, ihren Schnittpunkt gestrichelt mit der Spitze verbinden: das ist die Höhe, $2{,}4$ m; eine Grundkante mit $3{,}6$ m beschriften – nicht eine Seitenkante und nicht die Höhe einer Seitenfläche \quad n) M – in der Reihe D, N, M liegt genau eine Fläche dazwischen; N grenzt an D an}

9k)

\netzquader{2}{1}{0.5}{$4$ cm}{$2$ cm}{$1$ cm}


9l)

\quader{3}{2}{1.5}{$6$ cm}{$4$ cm}{$3$ cm}


9m)

\pyramide{3.5}{3.5}{3}{$3{,}6$ m}{}{$2{,}4$ m}

\erg{10}{a) Grund- und Deckfläche: Dreiecke (vorn und hinten); Seitenflächen: drei Rechtecke}
\erg{11}{a) $6$ cm und $3$ Kanten (Tiefe; gestrichelt)}
\erg{12}{a) Grundkreis auf dem Boden, er berührt alle vier Seitenwände (Durchmesser $50$ cm $=$ Kartonbreite), im Schrägbild als Ellipse; Spitze in der Mitte des Deckels (Schnitt der Diagonalen); zwei Mantellinien von der Spitze zum Kreisrand}
\erg{13}{a) Die zwei Dreiecke sind Grund- und Deckfläche, dazwischen liegen Rechtecke: Es ist ein Dreiecksprisma. Bei einer Pyramide laufen die dreieckigen Seitenflächen in einer Spitze zusammen. \quad b) Es hat nur fünf Quadrate; ein Würfel hat sechs Flächen, eine Seite bliebe offen. \quad c) $X = 5$, $Y = 4$, $Z = 6$ \quad d) Pyramide – ein Quadrat als Grundfläche, vier Dreiecke laufen in einer Spitze zusammen}
